\section{Four Dogmas within the One Mystery of Christ}

\dogmathesis{The Marian dogmas do not disclose an independent economy of salvation running beside Christ. They confess what the triune God has done in and for one created woman because the eternal Son truly became man from her, redeemed her by his merits, associated her freely and maternally with his mission, and brought her to the bodily glory promised to his Church. Every Marian ``yes'' is first God's gift; every Marian privilege manifests the sufficiency of Christ; every authentic Marian devotion ends in worship of Father, Son, and Holy Spirit.}

Catholic theology commonly identifies four Marian dogmas: Mary is truly Mother of God; she is ever-virgin; she was preserved from original sin from the first instant of her conception; and, when her earthly course was complete, she was assumed body and soul into heavenly glory. Calling all four \emph{dogmas} does not mean that each entered the Church's confession in the same way or received an identical kind of defining act. Divine motherhood was solemnly secured in the Church's Christological judgments at Ephesus and Chalcedon. Perpetual virginity belongs to ancient creedal, conciliar, liturgical, and universal teaching and received especially exact expression in the Lateran Synod of 649. The Immaculate Conception and Assumption were defined by solemn papal judgments in 1854 and 1950. Their unity is real; their histories and precise dogmatic objects remain distinct.

That distinction matters because a dogma is not simply a very cherished belief. It is a truth contained in divine Revelation and proposed by the Church as divinely revealed, to be believed with divine and Catholic faith (Second Vatican Council, \emph{Dei Verbum} [DV] 10; First Vatican Council, \emph{Dei Filius}, ch.~3; Congregation for the Doctrine of the Faith, 1998 doctrinal commentary on the concluding formula of the \emph{Professio fidei}, 5--11). The Church does not manufacture the reality by defining it. Under the assistance of the Holy Spirit, she identifies and binds the faithful to a truth already given in the apostolic deposit, sometimes by clarifying language made necessary through controversy and sometimes by bringing a long-received implication to definitive expression.

\subsection{Revelation, development, and definition}

Revelation is complete in the person and saving work of Jesus Christ (DV 2--4). Yet completion is not exhaustion of understanding. The apostolic Church received more than a list of future technical formulas: she received Christ, the Scriptures, sacramental life, apostolic preaching, and a rule of faith capable of faithful growth in understanding. The Church advances in perception of what has been handed on through contemplation and study, the spiritual understanding of believers, and the preaching of bishops in apostolic succession (DV 8). Authentic development preserves identity of truth while gaining precision. It neither adds a rival revelation nor excuses a historian from asking when a particular term or inference became explicit.

The Marian dogmas display several forms of development. The title \emph{Theotokos} protects the personal unity of the incarnate Word; its Marian force cannot be detached from the anti-Nestorian Christology that gave it conciliar precision. The threefold confession of virginity makes explicit what the Church perceived in the mystery of Christ's conception, birth, and Mary's consecrated vocation. The Immaculate Conception required centuries of dispute about original sin, the moment of animation, and the universality of Christ's redemption before preservative redemption supplied the decisive synthesis. The Assumption arose from the Church's worship, patristic meditation, and perception of Mary's unique relation to the risen Christ; its definition did not turn apocryphal narratives into history or invent a missing New Testament scene.

Development must therefore be narrated without two opposite errors. A flat antiquarianism says that unless a fourth-century Father uses the wording of 1854 or 1950, the doctrine is necessarily invented. A flat harmonization assigns the final formula to every ancient witness, hides disagreements, and makes controversy unintelligible. The responsible question is more exact: what reality does the witness attest; in what conceptual vocabulary; with what authority; and how does that witness belong to the later Church's judgment?

\subsection{The authority map used in this study}

\begin{enumerate}
  \item \textbf{Sacred Scripture} is inspired testimony and the soul of theology. A passage is read first in its literary setting, then within the canon. Typology is identified rather than disguised as literal narration.
  \item \textbf{A solemn definition or controlling magisterial act} states the dogmatic object. Its clauses and deliberate silences govern what may be attributed to the dogma.
  \item \textbf{The Fathers and liturgy} witness to received faith, patterns of interpretation, worship, controversy, and growth. Their authority is real but a selected catena is not a substitute for critical historical argument.
  \item \textbf{St.~Thomas Aquinas} supplies unsurpassed Christological, sacramental, and metaphysical principles. He is not made infallible by reverence, nor is his genuine disagreement with the later Immaculate Conception concealed.
  \item \textbf{Authoritative doctrines and titles} may be deeply rooted and binding without being additional Marian dogmas. Queenship, spiritual motherhood, lifelong personal holiness, and Mary's maternal intercession require their own classifications.
  \item \textbf{Theological synthesis, typology, and devotion} draw real conclusions and nourish prayer, but they remain accountable to the higher sources and cannot create a revealed fact.
\end{enumerate}

\subsection{Why the exact object matters}

A definition answers a delimited question. Ephesus does not say that Mary caused the divine nature. Perpetual virginity does not make marriage unholy or Mary less bodily. The Immaculate Conception does not mean that Mary was conceived virginally or saved without Christ. The Assumption does not define whether she died. Failure to preserve these limits turns dogma into caricature and makes apparent contradictions that the Church never asserted.

The same discipline governs Marian titles. A devotional term can have an orthodox intended meaning and still be judged pastorally or doctrinally unhelpful because its ordinary linguistic force obscures Christ. Historical papal usage is evidence of usage at a date; it does not by accumulation become a solemn definition. The final chapters accordingly receive Vatican II's account of Mary's maternal cooperation together with the Dicastery for the Doctrine of the Faith's 2025 doctrinal note \emph{Mater Populi Fidelis} [MPF], including its judgment that ``Co-redemptrix'' is always inappropriate as a title for expressing that cooperation (MPF 22).

\subsection{A source-audited, not self-authorizing, reference}

The prose identifies exact loci so that consequential claims can be checked. Short conventional translations convey the sense of sources whose original languages differ; the bibliography marks online patristic texts as working translations rather than critical editions. The accompanying research records disclose negative results, historical tensions, and unresolved critical-edition work. This study offers a Catholic theological synthesis governed by established doctrine. It does not claim ecclesiastical approval, substitute for the teaching office, or require assent to its original organization and analogies.
